%% ============================================================
%%  Chapter 4 — System Design (First Cut)
%% ============================================================
\chapter{System Design}
\label{ch:design}

\section{Design Overview}
\label{sec:design_overview}

\sys{} is designed as a four-phase, loosely-coupled pipeline running within a Kubernetes
cluster. The design adheres to the following principles:

\begin{description}
  \item[Kubernetes-nativeness:] All components are deployed as Kubernetes resources
  (DaemonSets, Deployments, ConfigMaps, CRDs). No external infrastructure is required.

  \item[Observability without modification:] Traffic telemetry is collected via eBPF
  without modifying application code, sidecar injection, or service mesh configuration.

  \item[Online operation:] All phases operate in a streaming fashion with bounded
  latency, enabling near-real-time detection and response.

  \item[Defence-in-depth:] Mitigation is graduated --- from soft throttling to selective
  blocking --- to minimise impact on legitimate traffic.

  \item[Reproducibility:] All components are open-source and deployable via standard
  Kubernetes manifests (Helm charts).
\end{description}

% ─────────────────────────────────────────────
\section{High-Level Architecture}
\label{sec:design_arch}

The system architecture is illustrated in Figure~\ref{fig:architecture}. The four phases
--- Observe, Model, Detect, Mitigate --- are described in detail in the following sections.

\begin{figure}[H]
\centering
% Replace with actual architecture diagram when available
% \includegraphics[width=\textwidth]{architecture}
\begin{tikzpicture}[
  node distance=0.8cm and 1.2cm,
  box/.style={rectangle, rounded corners=4pt, draw=NavyBlue, fill=NavyBlue!10,
              text width=3.2cm, minimum height=1.1cm, align=center, font=\small\bfseries},
  phase/.style={rectangle, rounded corners=2pt, draw=gray, fill=gray!10,
                text width=1cm, minimum height=4.5cm, align=center,
                font=\tiny\bfseries\color{gray}},
  arrow/.style={-Stealth, thick, NavyBlue},
  labelbox/.style={font=\tiny, text=gray}
]
  %% Workloads
  \node[box, fill=green!10, draw=ForestGreen] (pods) {Kubernetes\\Pods / Services};

  %% Phase 1 — Observe
  \node[box, right=1.5cm of pods] (ebpf) {eBPF Probes\\(Cilium/TC hooks)};
  \node[box, below=0.6cm of ebpf] (hubble) {Hubble\\Flow Exporter};
  \node[labelbox, above=0.1cm of ebpf] {Phase 1: Observe};

  %% Phase 2 — Model
  \node[box, right=1.5cm of ebpf, yshift=-0.5cm] (graphbuilder) {Graph\\Construction\\Engine};
  \node[labelbox, above=0.1cm of graphbuilder] {Phase 2: Model};

  %% Phase 3 — Detect
  \node[box, right=1.5cm of graphbuilder] (gnn) {GNN Anomaly\\Detector\\(GAE / GraphSAGE)};
  \node[labelbox, above=0.1cm of gnn] {Phase 3: Detect};

  %% Phase 4 — Mitigate
  \node[box, right=1.5cm of gnn, fill=red!10, draw=red!60] (mitigate) {Adaptive\\Mitigation\\Controller};
  \node[labelbox, above=0.1cm of mitigate] {Phase 4: Mitigate};

  %% Enforcement
  \node[box, below=1cm of mitigate, fill=orange!10, draw=orange!60] (xdp) {XDP/TC Rate\\Limiter +\\NetworkPolicy};

  %% Arrows
  \draw[arrow] (pods) -- (ebpf);
  \draw[arrow] (ebpf) -- (hubble);
  \draw[arrow] (hubble) -- (graphbuilder);
  \draw[arrow] (graphbuilder) -- (gnn);
  \draw[arrow] (gnn) -- node[above,font=\tiny]{scores} (mitigate);
  \draw[arrow] (mitigate) -- (xdp);
  \draw[arrow, dashed, red!50] (xdp.west) to[bend right=30] (pods.south);

\end{tikzpicture}
\caption{\sys{} high-level architecture: four-phase pipeline from eBPF telemetry to
adaptive mitigation.}
\label{fig:architecture}
\end{figure}

% ─────────────────────────────────────────────
\section{Phase 1: eBPF-Based Observability}
\label{sec:design_phase1}

\subsection{Telemetry Collection Strategy}

Rather than writing raw eBPF programs from scratch, this project leverages the Cilium CNI
plugin and its Hubble observability layer. This decision is justified by:

\begin{enumerate}
  \item \textbf{Production-readiness:} Cilium is a CNCF graduated project used in
  production by major cloud providers.
  \item \textbf{Rich flow metadata:} Hubble exports structured \texttt{flow.pb} protobuf
  records containing source/destination pod identities, protocol, verdict, L7 metadata.
  \item \textbf{gRPC streaming API:} Hubble's relay component exposes a gRPC streaming
  interface suitable for real-time consumption.
  \item \textbf{Reduced implementation risk:} Using Cilium's battle-tested eBPF programs
  frees development time for the novel graph ML and mitigation components.
\end{enumerate}

\subsection{Flow Record Schema}

Each Hubble flow record contains the fields shown in Table~\ref{tab:flow_schema}.

\begin{table}[H]
\centering
\caption{Key fields in a Hubble flow record used for graph construction}
\label{tab:flow_schema}
\begin{tabularx}{\textwidth}{llX}
\toprule
\textbf{Field} & \textbf{Type} & \textbf{Description} \\
\midrule
\texttt{source.pod\_name} & string & Source pod name \\
\texttt{source.namespace} & string & Source namespace \\
\texttt{destination.pod\_name} & string & Destination pod name \\
\texttt{destination.namespace} & string & Destination namespace \\
\texttt{l4.tcp/udp.destination\_port} & int & Destination port \\
\texttt{traffic\_direction} & enum & INGRESS / EGRESS \\
\texttt{verdict} & enum & FORWARDED / DROPPED / AUDIT \\
\texttt{summary} & string & L7 summary (HTTP method, DNS query) \\
\texttt{time} & timestamp & Flow event timestamp \\
\texttt{reply} & bool & Whether this is a reply flow \\
\bottomrule
\end{tabularx}
\end{table}

\subsection{Augmentation with Kubernetes Metrics}

Pod-level resource metrics (CPU usage, memory usage, network I/O) are collected from the
Kubernetes Metrics Server and Prometheus node-exporter, providing additional node features
for graph construction.

% ─────────────────────────────────────────────
\section{Phase 2: Dynamic Graph Construction}
\label{sec:design_phase2}

\subsection{Graph Topology}

The communication graph $G_t = (V_t, E_t, X_t, E\!F_t)$ is defined at each time window $t$ as:

\begin{itemize}
  \item $V_t$: Set of active pods and services observed in the window.
  \item $E_t$: Set of directed edges $(u, v)$ where pod $u$ communicated with pod $v$.
  \item $X_t \in \mathbb{R}^{|V_t| \times d_n}$: Node feature matrix.
  \item $EF_t \in \mathbb{R}^{|E_t| \times d_e}$: Edge feature matrix.
\end{itemize}

\subsection{Node Feature Vector ($d_n = 7$)}

\begin{table}[H]
\centering
\caption{Node (pod) feature vector components}
\label{tab:node_features}
\begin{tabular}{lll}
\toprule
\textbf{Feature} & \textbf{Source} & \textbf{Rationale} \\
\midrule
Request rate (req/s)        & Hubble flows      & Core volumetric indicator \\
Error rate (\%)             & Hubble verdicts   & Service degradation signal \\
Bytes received/s            & Hubble flows      & Bandwidth consumption \\
Bytes sent/s                & Hubble flows      & Amplification indicator \\
In-degree (active)          & Graph structure   & Number of incoming connections \\
Out-degree (active)         & Graph structure   & Number of outgoing connections \\
CPU utilisation (\%)        & Metrics Server    & Resource exhaustion indicator \\
\bottomrule
\end{tabular}
\end{table}

\subsection{Edge Feature Vector ($d_e = 5$)}

\begin{table}[H]
\centering
\caption{Edge (flow) feature vector components}
\label{tab:edge_features}
\begin{tabular}{lll}
\toprule
\textbf{Feature} & \textbf{Source} & \textbf{Rationale} \\
\midrule
Flow count (window)         & Hubble flows      & Connection frequency \\
Total bytes                 & Hubble flows      & Traffic volume \\
Protocol (one-hot)          & Hubble flows      & Protocol anomalies \\
Verdict ratio (drop/fwd)    & Hubble verdicts   & Policy violations \\
Packet size variance        & Hubble flows      & Unusual packet patterns \\
\bottomrule
\end{tabular}
\end{table}

\subsection{Sliding Window Strategy}

Graphs are constructed using a sliding window of $\Delta t = 30$ seconds with a step size
of $s = 10$ seconds (75\% overlap). This provides:
\begin{itemize}
  \item Temporal resolution sufficient for detecting rapid volumetric attacks.
  \item Sufficient data per window for graph construction with low-traffic microservices.
  \item Manageable graph update frequency for the GNN inference pipeline.
\end{itemize}

% ─────────────────────────────────────────────
\section{Phase 3: Graph-Based Anomaly Detection}
\label{sec:design_phase3}

\subsection{Detection Architecture}

The detection engine comprises two stages:

\begin{enumerate}
  \item \textbf{Stage 1 — Unsupervised Baseline (GAE):} A Graph Autoencoder trained
  exclusively on normal traffic snapshots. At inference time, high reconstruction error
  indicates anomalous graph structure or node features. This stage requires no labelled
  attack data.

  \item \textbf{Stage 2 — Supervised Classifier (GraphSAGE):} After generating labelled
  attack traffic in the testbed, a GraphSAGE node classifier is trained to distinguish
  attack types (volumetric, low-and-slow, application-layer). This stage provides
  fine-grained attack type identification.
\end{enumerate}

\subsection{Anomaly Score Computation}

For each graph snapshot $G_t$, the anomaly score for each node $v$ is computed as:

\begin{equation}
  s_v = \alpha \cdot \underbrace{\|\mathbf{x}_v - \hat{\mathbf{x}}_v\|_2}_{\text{feature reconst. error}}
       + (1-\alpha) \cdot \underbrace{\|\mathbf{A}_v - \hat{\mathbf{A}}_v\|_2}_{\text{structural reconst. error}}
\end{equation}

where $\hat{\mathbf{x}}_v$ and $\hat{\mathbf{A}}_v$ are the GAE reconstructions of
node features and adjacency, and $\alpha \in [0,1]$ balances feature vs. structural
anomaly weighting. The graph-level anomaly score is:

\begin{equation}
  S_G = \max_{v \in V} s_v
\end{equation}

\subsection{Threshold Calibration}

The detection threshold $\tau$ is set using the 99th percentile of anomaly scores on
held-out normal traffic: $\tau = \text{Percentile}_{99}(s_v \mid \text{normal})$.
An alert is raised when $s_v > \tau$ for any node $v$, with node identity preserved
for targeted mitigation.

% ─────────────────────────────────────────────
\section{Phase 4: Adaptive Mitigation Controller}
\label{sec:design_phase4}

\subsection{Tiered Response Policy}

The mitigation controller translates per-node anomaly scores into a tiered response:

\begin{table}[H]
\centering
\caption{Tiered mitigation response policy}
\label{tab:mitigation_tiers}
\begin{tabular}{cclc}
\toprule
\textbf{Tier} & \textbf{Score Range} & \textbf{Action} & \textbf{Enforcement} \\
\midrule
0 (Normal)     & $s_v < \tau$           & No action          & — \\
1 (Caution)    & $\tau \le s_v < 2\tau$ & Soft rate limit ($-25\%$ bandwidth) & TC eBPF \\
2 (Warning)    & $2\tau \le s_v < 3\tau$& Hard rate limit ($-75\%$ bandwidth) & TC eBPF \\
3 (Critical)   & $s_v \ge 3\tau$        & Selective block + alert operator    & XDP + NetworkPolicy \\
\bottomrule
\end{tabular}
\end{table}

\subsection{Enforcement Mechanisms}

\begin{description}
  \item[TC eBPF Rate Limiter:] A BPF program attached at the TC egress hook of each
  suspicious pod's veth interface implements a token-bucket rate limiter. The bucket
  refill rate is dynamically adjusted based on the anomaly score tier.

  \item[CiliumNetworkPolicy:] For Tier 3 alerts, the controller updates the
  \texttt{CiliumNetworkPolicy} for the affected pod, adding deny rules for the
  specific source IP or pod label responsible for the attack.

  \item[HPA Suppression Signal:] An annotation is added to the target deployment to
  suppress horizontal pod autoscaling during the attack, preventing cost amplification.
\end{description}

\subsection{Recovery Mechanism}

After $T_{recovery} = 60$ seconds of anomaly scores below $\tau$, the mitigation
controller removes rate-limiting rules and restores normal \texttt{CiliumNetworkPolicy}.
This prevents indefinite blocking of recovered services.

% ─────────────────────────────────────────────
\section{Deployment Architecture}
\label{sec:design_deployment}

\sys{} components are deployed as standard Kubernetes resources:

\begin{itemize}
  \item \textbf{Cilium + Hubble:} Installed via Helm as the cluster CNI.
  \item \textbf{Graph Builder:} Python \texttt{Deployment} consuming the Hubble gRPC stream.
  \item \textbf{GNN Detector:} Python \texttt{Deployment} (PyTorch + PyG) running inference
  on graph snapshots received from the Graph Builder via a shared message queue.
  \item \textbf{Mitigation Controller:} Python \texttt{Deployment} consuming anomaly scores
  and issuing Kubernetes API calls and eBPF program updates via \texttt{bpftool}.
  \item \textbf{Monitoring Stack:} Prometheus + Grafana \texttt{Deployment}s for metrics
  and dashboards.
\end{itemize}

% ─────────────────────────────────────────────
\section{Summary}
This chapter has presented the first-cut design of \sys{}, covering all four phases of the
pipeline with sufficient technical detail to guide implementation. Key design decisions
(use of Cilium/Hubble over raw eBPF, sliding-window graph construction, GAE+GraphSAGE
detection stages, tiered mitigation) have been justified with reference to the literature
review. Chapter~\ref{ch:implementation} will describe the implementation of each
component in detail.
